% ── Packages ──────────────────────────────────────────────────────────────────
\usepackage{xcolor}
\usepackage{tcolorbox}
\usepackage{amsmath}
\usepackage{amssymb}
\usepackage{amsthm}
\usepackage{mathtools}
\usepackage{enumitem}
\usepackage{algorithm}
\usepackage{algpseudocode}
\usepackage{listings}

\lstdefinestyle{pseudo}{
  basicstyle  = \ttfamily\small,
  mathescape  = true,
  columns     = flexible,
  keepspaces  = true,
  escapeinside= {(*}{*)},
  xleftmargin = 1.5em,
}

% ── Theorem environments ──────────────────────────────────────────────────────
\theoremstyle{plain}
\newtheorem{theorem}{Theorem}[section]
\newtheorem{lemma}[theorem]{Lemma}
\newtheorem{corollary}[theorem]{Corollary}
\newtheorem{proposition}[theorem]{Proposition}

\theoremstyle{definition}
\newtheorem{definition}{Definition}[section]
\newtheorem{example}{Example}[section]

\theoremstyle{remark}
\newtheorem*{remark}{Remark}
\newtheorem*{note}{Note}

% ── Colored boxes ─────────────────────────────────────────────────────────────

% \question{...}  – blue box for a question or concept to explore
\newcommand{\question}[1]{%
  \begin{tcolorbox}[
    colback  = blue!6!white,
    colframe = blue!45!black,
    boxrule  = 0.6pt,
    left = 6pt, right = 6pt,
    top  = 4pt, bottom = 4pt
  ]#1\end{tcolorbox}}

% \writing{...}  – plain indented block for working through an answer
\newcommand{\writing}[1]{%
  \begin{quote}\small #1\end{quote}}

% \followup{...}  – bullet list of follow-up questions to revisit
\newcommand{\followup}[1]{%
  \begin{itemize}[noitemsep, topsep=2pt]
    \item #1
  \end{itemize}}

% ── Common math shorthands ────────────────────────────────────────────────────
\newcommand{\R}{\mathbb{R}}
\newcommand{\C}{\mathbb{C}}
\newcommand{\N}{\mathbb{N}}
\newcommand{\Z}{\mathbb{Z}}
\newcommand{\Q}{\mathbb{Q}}
\newcommand{\F}{\mathbb{F}}
\newcommand{\set}[1]{\left\{ #1 \right\}}
\newcommand{\inner}[2]{\langle #1,\, #2 \rangle}

\DeclareMathOperator{\range}{range}
\DeclareMathOperator{\nullspace}{null}
\DeclareMathOperator{\rank}{rank}
\DeclareMathOperator{\kernel}{ker}
\DeclareMathOperator{\tr}{tr}
\DeclareMathOperator{\diag}{diag}
